\documentclass[10pt]{article}
\usepackage[margin=20mm]{geometry}
\usepackage{amsmath,amssymb,fontspec,unicode-math,xurl,hyperref}
\IfFontExistsTF{DejaVu Serif}{\setmainfont{DejaVu Serif}}{\setmainfont{Latin Modern Roman}}
\setmathfont{latinmodern-math.otf}
\setlength{\parindent}{0pt}\setlength{\parskip}{6pt}\sloppy
\title{Technical Verification of Feedforward Computational Graph Metrics}
\author{AI-assisted technical review}
\date{}
\begin{document}
\maketitle
Abstract: The two proposed graph metrics and the conditional polylogarithmic mixing argument were checked against their matrix definitions. Two local defects are demonstrated in the illustrative line-graph calculations: an incorrect mixing-time normalization and an off-by-one fidelity expression, together with an unattained maximum. The corrected calculations preserve the stated asymptotic contrast and do not invalidate the inspected FunSearch graph guarantee. The assessment is Minor; the empirical performance and implementation correspondence remain only partially verifiable.\par
\section*{Identity, scope and mathematical setting}
Alex Vitvitskyi, João Guilherme Madeira Araújo, Marc Lackenby and Petar Veličković, What Makes a Good Feedforward Computational Graph?, ICML 2025, PMLR 267, I25-43635 (0266), is reviewed in its 17-page proceedings version [paper]. Physical and printed page numbers agree. The independent check read the main argument and experimental description, pp. 1–9, and Appendices A–E, pp. 13–17. Original pages 4, 6, 8 and 13 were rendered to check the mathematical expressions, Figure 3 and proof layout. The frozen source manifest was read and its source identity preserved.\par
The graph has vertices \(0,\ldots,n-1\), a loop at every vertex, forward edges and a unique terminal vertex \(\tau=n-1\). With \(A_{ij}=1\) for \(j\to i\), the paper uses two different normalizations. The walk matrix \(W\) is column-stochastic, dividing by the source outdegree; the feature diffusion matrix \(\Delta\) is row-stochastic, dividing by the target indegree. Their roles cannot be interchanged. The stationary distribution for \(W\) is \(e_\tau\). Under the stated \(L^1\) convention, averaged mixing uses the first integer time at which \[L_n(t)=\frac1n\sum_i\|W^te_i-e_\tau\|_1<\frac14.\] The notation on pp. 3–4 sometimes calls the distance itself a mixing time; this report uses the threshold-time interpretation adopted by the subsequent arguments.\par
\section*{Two local line-graph corrections}
Problem 1 (The line-graph mixing calculation uses the wrong normalization: Minor).  Section 4.3.1, p. 4, derives an average non-absorption probability and concludes that the averaged mixing time is \(3n/2\).\par
Use the legal line graph with loops and edges \(i\to i+1\). Until absorption, a step advances with probability \(1/2\). If \(B_t\sim\operatorname{Bin}(t,1/2)\), averaging over the initial distance to \(\tau\) gives exactly \[S_n(t)=\frac1n\mathbb E[(n-1-B_t)_+]
=\frac{1}{n2^t}\sum_{j=0}^{\min(t,n-2)}(n-1-j)\binom tj .\] For any probability vector \(q\), \(\|q-e_\tau\|_1=2(1-q_\tau)\). Therefore the paper’s averaged distance is \(L_n(t)=2S_n(t)\), not \(S_n(t)\).\par
There is a separate range issue in the displayed simplification. The formula \(S_n(t)=1-(t+2)/(2n)\) is exact while \(t\le n-1\), since then the positive-part truncation changes nothing. It cannot be extrapolated as an exact identity to the claimed \(t\approx3n/2\). For larger \(t\), the omitted truncation correction is nonnegative.\par
The exact independent rational calculation in `check.py`/`check.json` yields:\par
\(n\) | Claimed \(3n/2\) | Correct threshold time | \(L_n(3n/2)\)\par
8 | 12 | 13 | \(2641/8192\)\par
16 | 24 | 27 | \(26200269/67108864\)\par
Both displayed distances exceed \(1/4\). At \(n=16\), the distances at times 26 and 27 are respectively \(19346383/67108864>1/4\) and \(8160303/33554432<1/4\). Every graph hypothesis is satisfied; this is not a disagreement about a disconnected graph or a different normalization.\par
The corrected leading coefficient also follows directly. For \(t/n\to c<2\), \(B_t/n\to c/2\) in probability and bounded convergence gives \(L_n(t)\to2-c\). Hence the threshold time divided by \(n\) tends to \(7/4\). A total-variation threshold would explain a leading coefficient \(3/2\), but p. 3 explicitly defines the unhalved \(L^1\) norm. That alternative therefore requires changing the convention.\par
The defect affects the exact calibration of an illustrative baseline. The line family’s linear-order mixing, the distinction from a rapidly mixing dense graph, and the conditional FS theorem do not depend on the printed coefficient. No dependency of the FS construction or empirical accuracy conclusions on that exact number has been established. This supports a local Minor assessment, rather than treating a correction to this example as a false central guarantee.\par
Problem 2 (The line-fidelity formula treats the source as a transient interior node: Minor).  Equation (6), p. 6, gives \(\max_a\binom{a}{n-1}/2^a\) for the line graph’s minimax fidelity. It conflicts with the row-normalized diffusion matrix.\par
For the same line, \(\Delta_{00}=1\), while \(\Delta_{ii}=\Delta_{i,i-1}=1/2\) for \(i\ge1\). The source is consequently different from the other nodes. For \(i\ge1\), writing \(d=n-1-i\) gives \[(\Delta^t)_{\tau i}=\binom td2^{-t}.\] By contrast, \((\Delta^t)_{\tau0}\) accumulates the mass reaching the absorbing source in the reversed walk and tends to one. It is not the transient coefficient with distance \(n-1\).\par
For \(n=3\), the exact matrix is \[\Delta=\begin{pmatrix}1&0&0\\[1pt]1/2&1/2&0\\[1pt]0&1/2&1/2\end{pmatrix}.\] The source-column supremum is one, the sink-column maximum is one at \(t=0\), and the middle-column maximum is \(1/2\) at \(t=1,2\). Thus the minimax value is \(1/2\), whereas Eq. (6) gives \(3/8\).\par
More generally, for \(n\ge3\) the worst column is \(i=1\), of distance \(n-2\). The exact corrected statement is \[\min_i\sup_{t\in\mathbb Z_{\ge0}}(\Delta^t)_{\tau i}
=\frac{\binom{2n-4}{n-2}}{4^{n-2}}.\] Indeed, for a positive integer \(d\), the ratio of consecutive coefficients is \((t+1)/(2(t+1-d))\); the maxima occur at \(t=2d-1,2d\). Their value \(\binom{2d}{d}/4^d\) decreases with \(d\), so the largest interior distance is worst. Direct exact matrix powers independently reproduce the formula for \(n=3,4,8,16\).\par
The replacement of max by sup is also necessary for the infinite-time definition: already for \(n=2\), the source coefficient is \(1-2^{-t}\), which never attains one. A declared finite horizon is another possible definition, but would be a different metric. The \(n=2\) line’s supremum-based minimax value is one, and is a separate boundary case.\par
Appendix B correctly maximizes a binomial coefficient with a fixed distance and applies Stirling’s formula. Its input distance is the part that does not match the graph’s worst column. Using \(n-2\) restores the exact calibration and still gives unnormalized fidelity asymptotic to \(1/\sqrt{\pi n}\) and normalized fidelity asymptotic to \(\sqrt{n/\pi}\). The paper’s qualitative line-versus-dense interpretation therefore survives. No released metric implementation was available to establish whether Figure 1 used Eq. (6), direct matrix powers or another horizon. A plotted-data error is not inferred from the algebra alone.\par
\section*{Other inspected theory and its boundaries}
\subsection*{Stationarity and loss of long-time fidelity}
Lemma 4.1’s stationary distribution proof is valid: any nonterminal mass at the smallest occupied vertex must decrease because its outdegree exceeds one. Proposition 5.1’s recurrence for the penultimate column gives a bound of the form \(t2^{-t}\), hence decay to zero, under its stated indegree assumption. These facts support examining transient fidelity rather than only its long-time limit.\par
The fully connected calculation on p. 4 contains a related local arithmetic clarification. With uniform initialization on the relabeled vertices \(0,\ldots,n\), the initial expected label is \(n/2\), and each step halves it; the expectation after \(t\) steps is \(n/2^{t+1}\). The printed \(n/2^t\) is still an upper bound, so the claimed logarithmic upper order remains valid. This is not evidence against fast mixing.\par
\subsection*{Path-count proposition}
The first-hit argument in Appendix A is sound for starts \(i<\tau\). Its averaging step silently includes the sink, where the bound needs an endpoint correction if \(s=1,\ldots,t\). Let \(P\) denote average absorption by \(t\). The weighted path sum is at least \(P-2^{-t}/n\), since the sink’s loop paths contribute \(1-2^{-t}\) rather than one. At the stated \(L^1\) mixing threshold, \(P>7/8\). For these nontrivial graphs the threshold cannot occur at \(t=1\), so \(t\ge2\) and \(2^{-t}/n\le1/8\). The sum still exceeds \(3/4\), and pigeonholing yields the stated factor \(3/(4t)\). Thus the endpoint omission does not refute Proposition 4.2.\par
\subsection*{Conditional FS mixing bound}
Theorem 6.1 assumes a fixed positive lower bound \(\alpha\) on the probability of crossing to the next block at each relevant level. Terminal blocks descend to the next recursive level rather than needing a nonexistent outward crossing. In the proof’s fixed-block-count model, with \(b=\lceil\log n\rceil\) blocks and \(D=O(\log n/\log\log n)\) levels, at most \(Db\) successful crossings are required. Conditional success probability at least \(\alpha\) gives \[\mathbb E H\le Db/\alpha,\qquad
L_n(t)\le2\Pr(H>t)\le \frac{2Db}{\alpha t}.\] Choosing \(t>8Db/\alpha\) proves the same polylogarithmic upper order and makes the \(L^1\) factor explicit. This argument concerns the pessimistic traversal bound; a lower bound on that comparison process does not automatically lower-bound the actual graph’s mixing time.\par
Remark C.1 explicitly acknowledges that equal edge allocations at all levels need not satisfy a constant \(\alpha\), and proposes a geometric allocation. The upper theorem remains conditional on the stated ratio and the construction details. It is not a theorem that every unqualified implementation or every sampled finite graph has that ratio. No substantive counterexample to the inspected conditional upper bound was established.\par
\section*{Empirical evidence and unresolved correspondence}
The inspected Figure 3 compares maximum retrieval, second-maximum retrieval and parity, with training lengths up to 256 and test lengths up to 1,024. Appendix E gives task-specific architectures, shared parameters for repeated maximum-retrieval blocks, an eight-layer parity model, optimizers and training budgets. Learning rate and weight decay were tuned on the fully connected graph and reused for the other masks. Those are relevant controls; the available text does not demonstrate test-set selection leakage.\par
The reported metric/performance alignment is empirical evidence, not a necessity or sufficiency theorem for learned computation. Exact edge budgets, random graph draws, seeds, metric horizon and raw score arrays were not available in a located official repository. The FS \(4\log n\) variant has a different nominal expander degree from a \(\log n\) variant, but the latter is also displayed. Without the recursive allocation and total-edge counts, an asserted unfair comparison would go beyond the evidence.\par
The Poisson generator accepts when a uniform draw exceeds \(p\), while some captions use \(0.2\) and Figure 3 uses \(0.8\). Acceptance versus rejection probability is an admissible explanation. The implemented convention remains Needs Evidence. The short Gemma experiment in Appendix E lacks the data needed for a separate reproducibility assessment. No model, training run or plotted curve was regenerated.\par
\section*{Assessment and state}
The highest supported assessment is Minor for two local analytical calibrations. The exact line formulas require correction, while the inspected structural propositions and conditional FS upper bound remain supported. Missing implementation details do not become empirical falsification claims.\par
Dimension | State\par
Issue assessment | Minor; two line-family calculation defects\par
Repair verification | Exact formulas derived; finite rational/matrix checks completed\par
Scientific decision | Pending; no human reviews\par
Author adoption | Not requested; unknown\par
Remaining evidence | Metric implementation, edge budgets, raw runs and historical source correspondence\par
9 A. Vitvitskyi, J. G. M. Araújo, M. Lackenby and P. Veličković. What Makes a Good Feedforward Computational Graph? ICML 2025, PMLR 267, 17-page proceedings PDF. https://proceedings.mlr.press/v267/vitvitskyi25a.html. OpenReview: `w9HPYVpfvY`.\par
\section*{Supporting checks}
The unchanged manifest and `source\_check.json` bind the staged original. Auxiliary `theory\_notes.json` and `independent\_notes.json` retain counterchecks and their limitations. No author code was executed.\par
Auxiliary source transport is mechanical. Bibliography originals were not exhaustively consulted. No official repository was identified in the acquired materials; absence of any release is not asserted. Separate attachments could not be enumerated after acquisition-time OpenReview HTTP 403. Later revisions and errata were not surveyed. No installation, new source download, author contact, publication or human adjudication occurred.\par
\end{document}
